\chapter{All testing charts for the unit-$\kappa,\nu,\ell$ source}
\label{component:p3f02c51074be}
\noindent\textbf{Component }\path{p3f02c51074be}\par
{\small\noindent Source: \path{NEW/c1/c144_unit_kappa_all_tests_rederived.tex}\par}
{\small\noindent SHA256: \path{3f02c51074beaad6ba4c99cea8fb710f31d0573f6bd28a7a49bb2457f076cbf7}\par}
{\small\noindent Inventory use: FU1--FU31 | FU6-FU23 exact finite-Bob parent/circuit identities, used only with freshly checked NU source limits | Fresh FU1-FU31 full unmatched and endpoint-testing proof; source endpoint restrictions retained\par}
\medskip
\section{All testing charts for the unit-$\kappa,\nu,\ell$ source}
\label{p3f02c51074be:sec:new-c144-unit-kappa-all-tests}

\paragraph{Exact scope and conclusion.}
Retain the original physical C14.4 source, its strict source gaps,
and its original prior-complement intersections. Assume
\[
 S=\sigma+L\in\operatorname{Inf}_{>0},\qquad
 \kappa=x/S,\quad\nu=\sigma/S,\quad\ell=L/S
 \quad\hbox{are positive bounded units},\qquad\nu+\ell=1.
 \tag{FU1}
\]
Together with the matched theorem UKR1--UKR21, the argument below
proves that the complete original parent is a finite separated tensor
SOS for every projective test and every state on FU1. It includes
arbitrary openings of $M,\theta,r,1-r$, literal $D=0$ and $\beta=0$,
and independent source and testing infinitesimals. It makes no claim
at the $\nu$- or $\ell$-infinitesimal source endpoints.

All normalizations are pure-source or pure-testing operations.
All real Schur complements below prove strictness of a real residue
matrix; no such inverse is used as a tensor-ring coefficient. The
strict bounded-Gram lemma is the one in the surviving
\path{inherited_r19/parent_reference.tex}.

\paragraph{The full parent and opening-safe source bounds.}
Write $D_u=1+u^2,D_v=1+v^2$, and divide the parent by $4$:
\[
\begin{split}
 F/4={}&m(vW-X+uvY-uZ)^2+CD_v(uW-Y)^2+AD_uD_vY^2\\
 &+(D+\tfrac12)(uW-Y+uvX-vZ)^2
 +(\alpha+\tfrac12)D_u(vY-Z)^2\\
 &+(\beta+\tfrac12)D_v(X+uZ)^2-\tfrac12D_uD_vZ^2.
\end{split}\tag{FU2}
\]
Keep the exact FM2--FM3 source formulas, with
$c_0=SC(1-r)$, $a_0=SAr$, $d_0=D/S$, and $a_1=\alpha/S$.
These four coefficients and $\beta\theta$ are bounded, with no
inverse of the possibly small $M$ or $\theta$ in their normalized
formulas. The $D/S$ numerator has its factor $S$ canceled exactly.

The following letters refer only to real residues:
\[
 A_* =\nu-\kappa h=\nu E,\quad R=1-A_*,\quad
 \nu=A_*+\kappa h,\quad\ell=1-A_*-\kappa h.
 \tag{FU3}
\]
The exact source yields
\[
\begin{gathered}
 E\ge1/18,\qquad h-A_*\ge0,\qquad
 A_*^2-\kappa(h-A_*)\ge0,\qquad R-\kappa\ge0,\\
 c_0=a_0=E/2,\qquad
 d_0=\{A_*^2-\kappa(h-A_*)\}/2,\\
 a_1=\frac{\ell\{\kappa h+(R-\kappa)A_*\}}{2\nu},
 \qquad 2\beta\theta=h-A_*.
\end{gathered}\tag{FU4}
\]
In particular $h\ge\nu/(1+\kappa)$, so $h,E,a_1,c_0,a_0$
have positive real lower bounds. The actual source has $m$ bounded
below by a positive real, while $C,A$ are unbounded. Here $m$
can be unbounded, and so can $\beta$. They cannot both be unbounded:
$m$ unbounded forces $R=\kappa$, whence
$\ell=\kappa\theta$ makes $\theta$ a unit; $\beta$ unbounded
forces $\theta=0$, whence $R-\kappa=\ell$ makes $M$ a unit.
For bounded $m,\beta$ their real values are
\[
 m=\frac{\kappa^2h}{2\nu(R-\kappa)},\qquad
 \beta=\frac{h-A_*}{2(1-h)}.
 \tag{FU5}
\]
These formulas are used only when the indicated real denominators
are positive. Every estimate below for the actual parent uses
bounded normalized source coefficients before taking residues.

\paragraph{A complete finite-Bob frame, including the state normals.}
Assume $u$ is infinitesimal and $v$ is finite. We may cap $m$ and
$\beta$ by finite pure-source values $m_c>0,\beta_c\ge0$ and
retain the complete excesses in their original event squares.
When a weight is bounded, keep it unchanged. When it is unbounded,
the sufficiently large real cap is chosen below. Define
\[
 B_0=X+uZ,\quad N_y=\frac{Y-ruW}{\sqrt{Sr(1-r)}},\qquad
 T=\kappa W,\ U=uZ,\ P=\frac{uW}{\sqrt S},\ Q=\sqrt S Z.
 \tag{FU6}
\]
All six rows are legal separated expressions, and $TU=\kappa PQ$
holds exactly. No restriction is placed on their values or on state
ratios. Set $a=u^2$, $s_y=\sqrt{Sr(1-r)}$, and
\[
\begin{split}
 M_0&=(v/\kappa)(1+ra)T-B_0,\\
 A_P&=D_v\{c_0(1-r)+a_0r(1+a)\}\\
 &\hspace{5mm}+S\{(D+\tfrac12)(1-r)^2
       +(\alpha+\tfrac12)(1+a)r^2v^2\}.
\end{split}\tag{FU7}
\]
The entire capped parent is $F_0+2L_yN_y+H_yN_y^2$, where
\[
\begin{split}
 F_0={}&m_cM_0^2+A_PP^2+(1+a)\{d_0(1+a)v^2+a_1\}Q^2
                   +\tfrac12(1+a)v^2U^2\\
 &-2(D+\tfrac12)(1+a)v^2B_0U
 -(1+a)v\{1+2D(1-r)+2\alpha r\}PQ\\
 &+\{(\beta_c+\tfrac12)D_v+(D+\tfrac12)av^2\}B_0^2
 +2(D+\tfrac12)(1-r)uv\sqrt S\,B_0P,
\end{split}\tag{FU8}
\]
\[
\begin{split}
 L_y={}&m_cuvs_yM_0
 +\sqrt{r(1-r)}D_v\{-c_0+a_0(1+a)\}P\\
 &+S\sqrt{r(1-r)}\{-(D+\tfrac12)(1-r)
          +(\alpha+\tfrac12)(1+a)rv^2\}P\\
 &-(D+\tfrac12)uvs_yB_0
       +\sqrt{r(1-r)}v(1+a)(D-\alpha)Q,\\
 H_y={}&rc_0D_v+(1-r)a_0(1+a)D_v\\
 &+Sr(1-r)\{m_cu^2v^2+D+\tfrac12
                      +(\alpha+\tfrac12)(1+a)v^2\}.
\end{split}\tag{FU9}
\]
These are exact identities in the two unrestricted original state
coordinates $X,Y$; $B_0$ is not set to zero. Their coefficients are
bounded across all source openings. In $L_y$ every coefficient has
zero residue, since $c_0,a_0$ have the same residue, and the other
terms contain $u$ or $S$. The $N_y$ pivot has residue $ED_v/2>0$.
Thus it suffices to make the remaining five-row real residue strict.

\paragraph{Finite nonzero unmatched Bob residue.}
Take $v>0$ with a positive finite residue different from
$b=\sqrt{\ell/\nu}$. The complete real residue of FU8--FU9 is
\[
\begin{split}
 &m_c\{(v/\kappa)T-B_0\}^2+qD_vB_0^2-v^2B_0U
       +\tfrac12v^2U^2-vPQ\\
 &\hspace{8mm}+\tfrac E2D_vP^2+(d_0v^2+a_1)Q^2
             +\tfrac E2D_vN_y^2,\qquad q=\beta_c+\tfrac12.
\end{split}\tag{FU10}
\]
For the moment use the uncapped finite values $m,\beta$ from FU5.
Eliminating $B_0$ in this real matrix, and using $TU=\kappa PQ$
only to name the adjusted cross coefficient, gives
\[
 a_TT^2+a_UU^2+a_PP^2+a_QQ^2-2fPQ,
 \tag{FU11}
\]
where, for $c=v^2$ and $H=m+q(1+c)$,
\[
\begin{gathered}
 a_T=\frac{mq(1+c)c}{\kappa^2H},\quad
 a_U=\frac{c(2H-c)}{4H},\quad a_P=E(1+c)/2,\quad
 a_Q=d_0c+a_1,\\
 f=\frac{v}{2}(1+mc/H).
\end{gathered}\tag{FU12}
\]
The replacement of the real Schur cross term $TU$ by $\kappa PQ$
is implemented in the actual bounded Gram by a real multiple of
the exact zero $TU-\kappa PQ$. The elimination of $B_0$ itself
is only a real strictness test. Every diagonal is positive,
including when $D$ or $\beta$ is zero.
The reduced determinant in a formal positive real variable $y$ is
\[
 \Delta(y,c)=(a_T\kappa^2+a_Py)(a_Uy+a_Q)-f^2y.
 \tag{FU13}
\]
It obeys $4H\Delta/(1+c)=\mathcal P(y,c)$ with
\[
 \mathcal P=p_{22}y^2c^2+p_{21}y^2c+p_{12}yc^2
              +p_{11}yc+p_{10}y+p_{02}c^2+p_{01}c,
 \tag{FU14}
\]
\[
\begin{array}{ll}
 p_{22}=E\beta,&p_{21}=E(m+q),\\
 p_{12}=2m\beta+2qEd_0,&
 p_{11}=-m-q+2E\{(m+q)d_0+qa_1\},\\
 p_{10}=2(m+q)Ea_1,&p_{02}=4mqd_0,\quad p_{01}=4mqa_1.
\end{array}\tag{FU15}
\]
All coefficients except possibly $p_{11}$ are nonnegative, and
$p_{21},p_{10},p_{01}$ are strictly positive. Direct substitution of
FU3--FU5 gives the exact three contact identities
\[
 \mathcal P(\kappa,\ell/\nu)=
 \mathcal P_y(\kappa,\ell/\nu)=
 \mathcal P_c(\kappa,\ell/\nu)=0.
 \tag{FU16}
\]
These finite rational identities can also be checked without
radicals. For positive real $y,c$,
\[
 \frac{\mathcal P(y,c)}{yc}
 =p_{22}yc+p_{21}y+p_{12}c+p_{11}
                     +p_{10}/c+p_{02}c/y+p_{01}/y.
 \tag{FU17}
\]
As a real function of $(\log y,\log c)$, this is strictly convex:
each nonconstant summand has nonnegative Hessian, and
$p_{21}y+p_{10}/c$ alone has positive definite Hessian. FU16 says
its value and gradient vanish at the contact. It is consequently
positive at every other positive point. This real argument is only
a way to determine a finite residue matrix, not a use of logarithms
or mixed inverses in the tensor ring.

Since $c\ne\ell/\nu$, FU13 is positive for all $y>0$.
Minimizing its quotient by $y$ gives
\[
 f<\kappa\sqrt{a_Ta_U}+\sqrt{a_Pa_Q}.
 \tag{FU18}
\]
Indeed the minimum of
$a_Pa_Uy+a_T\kappa^2a_Q/y$ is
$2\kappa\sqrt{a_Ta_Ua_Pa_Q}$, attained at a positive $y$;
equality in FU18 would give a zero of FU13. Choose a real
$\lambda$ with $|\lambda|<\sqrt{a_Ta_U}$ and
$|f-\kappa\lambda|<\sqrt{a_Pa_Q}$. Adding the exact zero
$-2\lambda(TU-\kappa PQ)$ makes both $2\times2$ blocks in
FU11 strictly positive. Restoring the positive real $B_0$ pivot
and the $N_y$ pivot proves strictness of the complete six-row
residue. All actual Gram entries were bounded in FU8--FU9, so
the strict bounded-Gram lemma gives the tensor SOS.

\paragraph{Unbounded $m$ or $\beta$ in the unmatched chart.}
The preceding strictness persists under a sufficiently large
finite cap in the one possible unbounded weight. This assertion
has an opening-safe direct verification. If $m\to\infty$,
divide FU14--FU15 by $m$ before taking the real limit. Then
\[
 (p_{22},p_{21},p_{12},p_{11},p_{10},p_{02},p_{01})/m
 \longrightarrow
 (0,E,2\beta,-1+2Ed_0,2Ea_1,4qd_0,4qa_1).
 \tag{FU19}
\]
At $R=\kappa$, the three contact identities FU16 remain exact
for this limit polynomial. The strict convexity proof FU17 still
applies because its $y$ and $1/c$ coefficients are positive.
The Schur coefficients tend to
\[
 a_T\to q(1+c)c/\kappa^2,\qquad
 a_U\to c/2,\qquad f\to v(1+c)/2,
 \tag{FU20}
\]
with $a_P,a_Q$ unchanged. Therefore FU18 is strict at the limit,
and continuity makes it strict for a sufficiently large finite
real $m_c$.

If $\beta\to\infty$, divide instead by $q$ before taking the
limit at $h=1$. The corresponding coefficient vector is
\[
 (E,E,2m+2Ed_0,-1+2E(d_0+a_1),
                         2Ea_1,4md_0,4ma_1),
 \tag{FU21}
\]
again with the three exact contact identities and the two strictly
positive convexity coefficients. Now
$a_T\to mc/\kappa^2$, $a_U\to c/2$, $f\to v/2$.
A sufficiently large finite $\beta_c$ therefore works.
All the other normalized actual coefficients have their real
residues with bounded unit denominators; no $O(S)$ term was
divided by $M$ or $\theta$. Restore the complete excesses
$(m-m_c)(vW-X+uvY-uZ)^2$ and
$(\beta-\beta_c)D_v(X+uZ)^2$. These caps are used only on the
unmatched chart and need not vanish at the actual contact.

\paragraph{Infinitesimal or zero Bob coordinate.}
If $v$ is infinitesimal, use FU8--FU9 with the six rows
\[
 (B_0,T_v,U_v,P,Q,N_y),\qquad T_v=vT,\quad U_v=vU.
 \tag{FU22}
\]
Cap $m$ at any sufficiently small fixed positive real $m_c\le m$,
and discard no excess. One may set $\beta_c=0$ and retain its
whole event. The exact coefficients remain bounded: the $B_0U$
coefficient is now a multiple of $vB_0U_v$, and $M_0$ becomes
$(1+ra)T_v/\kappa-B_0$. Their real residue is
\[
 m_c(T_v/\kappa-B_0)^2+\tfrac12B_0^2+\tfrac12U_v^2
                      +\tfrac E2(P^2+N_y^2)+a_1Q^2>0.
 \tag{FU23}
\]
This also covers literal $v=0$. The rows need not be independent
coordinates on the state space.

\paragraph{Infinitesimal reciprocal Bob coordinate.}
Put $w=1/v$ and use the homogeneous parent $w^2F$ at $w=0$ as
well as for nonzero infinitesimal $w$. The complete signed core
in the quarter-parent is $K_w/2$, with
\[
\begin{split}
 K_w={}&u^2H^2-2u^2HY-2uHW+2uwHX+u^2w^2W^2\\
 &+2u^2wWX-2uw^2WY+(1+u^2+w^2)(X^2+Y^2)-2uwXY,
 \qquad H=wZ.
\end{split}\tag{FU24}
\]
This is the exact CE6 full core, so all state normals are present.
FU4 gives the real alternative
\[
 d_0+\kappa\beta\theta=(\nu E)^2/2>0.
 \tag{FU25}
\]
Thus either $D/S$ has a positive real lower bound, or $\beta$
has one. Both cases allow independent rates for $u,w,S$.

If $\beta$ is strong, use the four rows $(W,X,Y,U)$, $U=uZ$,
in FU24 by replacing $uH$ with $wU$. Every coefficient is bounded
and the core residue is $X^2+Y^2$. The exact $m$ event is
$(W-wX+uY-wU)^2$, and the $\beta$ event is
$D_w(X+U)^2$. Positive finite caps give the strict residue
\[
 m_cW^2+\tfrac12(X^2+Y^2)+\beta_c(X+U)^2>0.
 \tag{FU26}
\]
All other event packets and both excesses remain whole.

If $D/S$ is strong, choose positive real $c,d,m_c$ with
$c\le c_0,a_0$, $d\le D/S$, and $m_c\le m$. Select exactly
\[
 C_c=c/[S(1-r)],\qquad A_c=c/(Sr),\qquad D_c=dS.
 \tag{FU27}
\]
Retain all excesses, including the whole $A_cu^2D_wY^2$.
Writing $Y=ruW+\sqrt{Sr(1-r)}N_y$, the other two selected
$C,A$ terms have the exact harmonic identity
\[
 C_cD_w(uW-Y)^2+A_cD_wY^2=cD_w(P^2+N_y^2).
 \tag{FU28}
\]
Use six rows $(W,V,X,P,N_y,Q)$ with $V=uwZ$ and the exact
monomial identity $VW=wPQ$. Replacing $uH$ by $V$ in FU24,
the potentially non-small cross term $-2VW$ is exactly $-2wPQ$.
Every other off-diagonal core coefficient is infinitesimal after
$Y=\sqrt S\{rP+\sqrt{r(1-r)}N_y\}$; its residue is $V^2+X^2$.
The selected $m$ and $D$ events are the complete squares
\[
 m_c(W-wX+uY-V)^2,\qquad
 d\{uw\sqrt S W+u\sqrt S X-w\sqrt S Y-Q\}^2.
 \tag{FU29}
\]
The full selected residue is strictly positive:
\[
 m_c(W-V)^2+\tfrac12(V^2+X^2)+c(P^2+N_y^2)+dQ^2>0.
 \tag{FU30}
\]
The strict bounded-Gram lemma applies to the exact complete
forms, and all unselected positive packets are restored. Both
reciprocal proofs include literal $w=0$ and $u=0$.

\paragraph{Exhaustion of tests and exact joins.}
If $u$ is noninfinitesimal, $C,A$ are unbounded by FU4. The
surviving two-adjacent-events theorem CE1--CE10 in
\path{common_remote.tex}
therefore applies to the original full parent and all states.
Its hypotheses require only nonnegative event weights and these
two unbounded coefficients; no small-$\kappa$ source statement
is imported.

For infinitesimal $u$, every positive projective $v$ is either
infinitesimal, a finite unit, or has infinitesimal reciprocal.
FU22--FU30 cover the two end cases. A finite unit has either the
matched residue, covered by UKR1--UKR21, or a different residue,
covered by FU10--FU21. The transformations
\[
 (u,Y,Z)\mapsto(-u,-Y,-Z),\qquad
 (v,X,Z)\mapsto(-v,-X,-Z)
 \tag{FU31}
\]
preserve FU2 exactly and handle both signs. Literal axes and
projective endpoints were included in the displayed raw frames.
Complex states split into real and imaginary parts. No proof in
this exhaustion bounds, normalizes, or relates the testing-state
ratios $W/Z$, $X/Z$, or $Y/Z$.

Consequently every testing chart is covered on exactly FU1,
with every original source predicate retained. The remaining
source endpoints are separate from this conclusion.
